% Generated by report/make_restated.py from theory/B2_effective_curvature.tex: do not edit (edit the note and rerun).
\begin{restated}{Remark}{rem:B2-what-explains}{B2}{0}{what each constant can explain}
With the step size fixed at $\gamma=1/(4R^2)$, the gain coming from $\LEff<R^2$ is the factor
$(8R^2-5\LEff)/(3R^2)\in[1,8/3]$ on the parameter-term exponent: the co-coercivity slack can explain at most a
factor $8/3$ of the gap, however small $\LEff$ is (a smaller $\LEff$ would allow a larger step, but the
algorithm does not take it). In contrast, the gain from $\muEff>\mu$ is the factor $\muEff/\mu$, which can be as
large as $L/\mu$ (Lemma~\ref{lem:B2-mueff}) and is typically of the order described in
Sections~\ref{sec:B2-logistic}--\ref{sec:B2-direction}. The parameter term is the active one in the minimum as
long as $a(\theta_{t-1})<1/(3n)$, i.e.\ (with $\LEff=R^2$) $\muEff(\theta_{t-1})<16R^2/(9n)$; in the regime
$\kappa\gg n$ of hypothesis (A4) this holds for $\mu$ and for moderate values of $\muEff$, and the bound saturates
at $1-\frac1{3n}$ once $\muEff$ exceeds $16R^2/(9n)$: beyond that point no improvement of the curvature term can be
``explained'', and the residual gap belongs to the memory coefficient $1-\frac1{3n}$, which is set by the square-norm
inequality \eqref{ineq:B3-I1}, the inequality $\Ei\norm{Z-\Ei Z}^2\le\Ei\norm Z^2$ \eqref{ineq:B3-I3} and the choice of
$\Delta$, and to the min step \eqref{ineq:B3-I5} between the two rates, all instrumented in the B3 note. In its
experiments with $\kappa\le10^2$, where the memory term is active from the start, order~A of
Proposition~\ref{prop:B3-attribution} (which, as here, keeps the minimum until the end) assigns $0.56$--$0.61$ of the
logarithmic gap to (I5), $0.38$--$0.44$ to (I1) and almost nothing to (I3).
\end{restated}
